In this appendix we reformulate the $r$-matrix structure in operator form \cite{SemenovTyanShanskii1983,babelon2003introduction}.

We fix a matrix realization $\mathcal G \subset \mathfrak{gl}(r+1)$ and use the invariant bilinear form
\begin{equation}
    (X,Y)=\Tr(XY),
\end{equation}
which identifies $\mathcal{G} \simeq \mathcal{G}^*$.

\begin{definition}
    For $L \in \mathcal{G}$ and $X \in \mathcal{G}$, we define the scalar function
    \begin{equation}
        L(X) \equiv (L,X).
    \end{equation}
\end{definition}

Using this identification, any element $r_{12} \in \mathcal{G}^{\otimes 2}$ defines a linear operator
$R:\mathcal{G} \to \mathcal{G}$ by
\begin{equation}
    R(X)=\Tr_2\left(r_{12}(1 \otimes X)\right).
\end{equation}
This provides an equivalent operator formulation of the $r$-matrix formalism.

\begin{proposition}
    Let $r_{12} \in \mathcal{G}^{\otimes 2}$ and let $R$ be the associated operator.  
    If the Poisson brackets of $L \in \mathcal{G}$ satisfy
    \begin{equation}
        \{L_1,L_2\} = [r_{12},L_1] - [r_{21},L_2],
    \end{equation}
    then, for all $X,Y \in \mathcal{G}$,
    \begin{equation}
        \{L(X),L(Y)\} = L([X,Y]_R),
    \end{equation}
    where
    \begin{equation}
        [X,Y]_R = [R(X),Y] + [X,R(Y)].
    \end{equation}
\end{proposition}

\begin{proof}[Sketch of proof]
    Pairing the tensorial identity with $X \otimes Y$ and using the invariance of the scalar product,
    one rewrites both sides in terms of the pairing $(L,X)$. The commutator terms reduce to expressions involving $R$, yielding the stated formula.
\end{proof}

The Yang--Baxter equation admits a natural formulation in terms of the operator $R$.

\begin{proposition}
    Let $r_{12} \in \mathcal{G}^{\otimes 2}$ be antisymmetric and let $R$ be defined as above.  
    If $r_{12}$ satisfies the modified classical Yang--Baxter equation
    \begin{equation}
        [[r,r]] = \mu \Omega,
    \end{equation}
    then $R$ satisfies
    \begin{equation}
        [R(X),R(Y)] - R\left([R(X),Y] + [X,R(Y)]\right) = -\mu [X,Y]
        \qquad \forall X,Y \in \mathcal{G}.
    \end{equation}
\end{proposition}

\begin{proof}[Sketch of proof]
    Pairing the tensorial identity $[[r,r]]=\mu\Omega$ with $X \otimes Y \otimes Z$
    and using invariance, one translates the tensorial commutators into operator expressions involving $R$.
\end{proof}

In the case of the Toda chain, the $r$-matrix constructed in Section~\ref{sec:toda_rmatrix}
defines an operator $R$ diagonal in the root space decomposition:
\begin{equation}
    R(H)=0, \qquad
    R(E_{\pm\beta}) = \pm \frac12 E_{\pm\beta}
    \qquad (\beta \text{ positive}).
\end{equation}
This provides an explicit solution of the modified classical Yang--Baxter equation. 